\section{The local annihilator when zero is absent from one set of jumps}

Write $p$ for the local prime, $K_0$ for the maximal unramified subfield
of $K$, $K_j=K_0(\mu_{p^j})$, $K_\infty=\bigcup_j K_j$,
$G_j=G_{K_j}$ and $H=G_{K_\infty}$. Absolute values satisfy $|p|=p^{-1}$.
All unit balls below are coordinatewise balls for this absolute value.

% BEGIN A79L1
\paragraph{A79L1. Integral averaging over the cyclotomic tower.}
If $L/K_\infty$ is finite Galois, there exists $b\in\mathcal O_L$
with $\operatorname{Tr}_{L/K_\infty}(b)=p$.
We give the ramification-theoretic proof of the particular trace value
needed here. Take a finite Galois extension $F/K_0$ in the algebraic
closure containing generators of $L/K_\infty$. Put $E_j=FK_j$. It is
enough to find the trace value $p$ in $FK_\infty$, since trace is
transitive and preserves integrality. For large $j$ the decreasing
groups $\operatorname{Gal}(E_j/K_j)$, embedded in
$\operatorname{Gal}(F/K_0)$ by restriction, stabilize at
$\operatorname{Gal}(FK_\infty/K_\infty)$. Indeed the fields
$F\cap K_j$ stabilize at $F\cap K_\infty$.

Choose a positive integer $a$ beyond the last upper ramification break
of $F/K_0$. The different exponents of $E_j/K_j$, measured with
the integer valuation of $E_j$, are bounded independently of $j$.
Here are the numerical details. Write $G=\operatorname{Gal}(E_j/K_0)$,
$H_j=\operatorname{Gal}(E_j/K_j)$ and $\Gamma_j=G/H_j$.
The quotient theorem for upper ramification groups gives
$G^a\cap H_j=1$, since it acts trivially on both $F$ and $K_j$.
Herbrand's subgroup formula gives
$H_j^{\psi_{\Gamma_j}(a)}=H_j\cap G^a$. The numbers
$\psi_{\Gamma_j}(a)=\int_0^a[\Gamma_j^0:\Gamma_j^v],dv$
are bounded independently of $j$. In fact a cyclotomic automorphism
of exponent $c$ with $v_p(c-1)=k<j$ moves the uniformizer
$\zeta_{p^j}-1$ by an element of integer valuation $p^k$.
Thus the lower breaks occur at $p^k-1$; integrating the lower
filtration gives upper breaks at integers and
$[\Gamma_j^0:\Gamma_j^v]\le p^{\lceil a\rceil+1}$ on $[0,a]$.
This calculation also allows the initial tame factor of size $p-1$.
If $D=[F:K_0]$, then $|H_j|\le D$, so passing from upper to lower
indices for $H_j$ multiplies the bound on the last break by at most
$D$. Hilbert's different formula $\sum_{n\ge0}(|(H_j)_n|-1)$
now gives a uniform integer bound $C$ on the different exponent.

For a finite extension of local fields $E/F'$ of ramification
index $e$ and different exponent $d$, trace duality gives
$\operatorname{Tr}(\mathcal O_E)=\mathfrak p_{F'}^{\lfloor d/e\rfloor}$.
Indeed containment in $\mathfrak p_{F'}^t$ is equivalent to
$\mathfrak p_{F'}^{-t}\mathcal O_E\subseteq\mathfrak D^{-1}_{E/F'}$.
Here $v_{K_j}(p)=\varphi(p^j)$ eventually exceeds $C$, so
$p\in\operatorname{Tr}_{E_j/K_j}(\mathcal O_{E_j})$. At a stabilized
level the trace sum has exactly the same conjugates upstairs.
This proves the assertion.
% END A79L1

% BEGIN A79L2
\paragraph{A79L2. Descent of a vector fixed modulo the unit ball.}
Set $C=\mathbb C_p$, $B_r=\mathcal O_C^r$. Suppose that
$x\in C^r$ satisfies $hx-x\in B_r$ for every $h\in H$, with
coordinatewise action. Then $px=z+e$ with $z\in K_\infty^r$
and $e\in B_r$. Choose an algebraic vector $x'$ at distance
at most one from $x$, and a finite Galois $L/K_\infty$ containing
its coordinates. Take $b$ from A79L1 and put
$z=\operatorname{Tr}_{L/K_\infty}(bx')$ coordinatewise. For each
conjugate $s$, lift $s$ to $H$. Then $|s(x')-x|\le1$ in every
coordinate, since $|s(x'-x)|\le1$ and $|s(x)-x|\le1$.
Consequently $px-z=\sum_s s(b)(x-s(x'))\in B_r$.
The same argument with $hx=x$ and approximation error $\epsilon$
bounds $|px-z|$ by $\epsilon$. It follows that
$C^H=\widehat K_\infty$, the closure inside $C$; the reverse inclusion
follows from continuity. Only these two scalar descent statements
are needed for the untwisted invariant classes.
% END A79L2

% BEGIN A79L3
\paragraph{A79L3. Inverting the generator on the complementary summand.}
Let $m\ge1$, with $m\ge2$ at $p=2$. Choose $g\in G_m$ for
which $v_p(\chi(g)-1)=m$, and put $\delta=\zeta_p-1$.
Write $\Gamma$ for the action of $g$ on the tower. There is an integral
splitting
\[
 \mathcal O_{\widehat K_\infty}=\mathcal O_{K_m}\oplus X .
\]
The projection at level $n\ge m$ is normalized trace to $K_m$.
Indeed each step $K_j/K_{j-1}$, $j>m$, has integral basis
$1,\zeta_{p^j},\ldots,\zeta_{p^j}^{p-1}$. To see integrality,
$\zeta_{p^j}-1$ is a uniformizer, and its first $p$ powers have distinct
valuations modulo $p$; replacing these by the powers of $\zeta_{p^j}$
is a triangular integral change of basis. The normalized trace
retains the constant coefficient. These integral projections commute
with $g$ and with one another. They are uniformly continuous in the
$p$-adic topology since they preserve $p^t$ times the unit ball for
each $t$. Thus the splitting extends to the completion.

On $X$ the equation $(\Gamma-1)y=\delta w$ has a unique solution
$y\in X$ for every $w\in X$. First work at a finite level and split
the kernel into its successive trace kernels. On the piece
$\zeta_{p^j}^b K_{j-1}$, $1\le b<p$, put $h=p^{j-m-1}$.
The binomial valuation formula gives
$v_p(\chi(g)^h-1)=j-1$, so $g^h$ fixes $K_{j-1}$ and acts on
the piece by a primitive $p$-th root $\eta$. An inverse to
$\Gamma-1$ there is $(1+\Gamma+\cdots+\Gamma^{h-1})/(\eta-1)$.
This formula applies to the piece since $g$ preserves it
($\chi(g)\equiv1\pmod p$). As $\delta/(\eta-1)$ is a unit of
$K_m$, the inverses times $\delta$ preserve the integral trace kernel.
They agree under level inclusion by uniqueness, and extend to $X$
by uniform continuity. Their two inverse identities extend by density.

Let $A$ be a square matrix over $\mathcal O_{K_m}$ with
$|A-I|_{\max}<|\delta|=p^{-1/(p-1)}$. On $X^r[1/p]$,
$A\Gamma-I$ is bijective and $\delta(A\Gamma-I)^{-1}$ preserves
$X^r$. In fact
\[
 A\Gamma-I=(I+(A-I)\Gamma(\Gamma-I)^{-1})(\Gamma-I).
\]
The entries of $(A-I)/\delta$ lie in the maximal ideal of
$\mathcal O_{K_m}$. Hence the correction in parentheses differs
from $I$ by an operator sending $X^r$ into $\varpi_m X^r$, where
$\varpi_m$ is a uniformizer of $K_m$. The inverse geometric series
converges in this complete module. The same conclusion holds for
$\chi(g)^u A$ with any $u\in\mathbb Z$, since
$|\chi(g)^u-1|\le p^{-m}<|\delta|$ at the permitted levels.
% END A79L3

% BEGIN A79L4
\paragraph{A79L4. Distinct weights control the finite-level denominator.}
Suppose $U:G_m\to\operatorname{GL}_r(\mathcal O_{K_m})$ is a
representation trivial on $H$, with $|U(g)-I|_{\max}<|\delta|$.
Put $A=U(g)$ and $D_g=A\Gamma$. Suppose $C^r$ has a basis
$e_k$ of weighted vectors, with
$\chi(\sigma)^{u_k}U(\sigma)\sigma(e_k)=e_k$ on $G_m$, and
let $J$ contain the distinct weights. Assume $J$ is finite and
$0\notin J$. Each $e_k$ belongs to $\widehat K_\infty^r$ by A79L2.
Split it into its $K_m$ and complementary parts using A79L3 after
clearing denominators. The latter part is killed by
$\chi(g)^{u_k}A\Gamma-I$, so it vanishes. Thus
$Ae_k=\chi(g)^{-u_k}e_k$. In particular
\[
 F(A)=0,\qquad F(Y)=\prod_{u\in J}(Y-\chi(g)^{-u}).
\]
Put $N=\sum_{u\in J}(m+v_p(u))$. For each $u\ne0$,
$v_p(1-\chi(g)^{-u})=m+v_p(u)$. Raising $1+p^m t$, $t$ a unit,
to an exponent prime to $p$ preserves the valuation of its difference
from $1$. Raising to $p$ increases that valuation by one, by the
binomial theorem at the stated levels. Inversion preserves it as well.
Hence $v_p(F(1))=N$. For $r>0$, $J$ is nonempty and $N\ge1$.
On the finite component the identity
$(A-I)((F(Y)-F(1))/(Y-1))(A)=-F(1)I$ shows that
$p^N(A-I)^{-1}$ is integral. On the complement A79L3 gives the
same assertion for $D_g-I$ because $p^N/\delta$ is integral.
It follows that whenever $z\in\widehat K_\infty^r$ satisfies
$(D_g-I)z\in\mathcal O_{\widehat K_\infty}^r$, then
$p^Nz\in\mathcal O_{\widehat K_\infty}^r$. Repeated weights cost
only one factor in $F$.
% END A79L4

% BEGIN A79L5
\paragraph{A79L5. The exponent on the original rational lattice.}
Use a $\mathbb Z_p$-basis of the original full lattice $T$ and its
$\mathbb Q_p$-basis on $V=T[1/p]$; write $M(\sigma)$ for the
action matrix and $d$ for the rational dimension.
Fix an embedding $\tau:K\to L$ and its finite set $U_\tau$ of jumps
with $0\notin U_\tau$. Use the following precise component comparison
at $m=m_N(p,s,M_0+h_0)$ with $K\subset K_m$.
There are $r,P:C^d\to C^r,g,U,J$ as in A79L3--A79L4,
$J\subseteq U_\tau$, and
\[
 P M(\sigma)=U(\sigma)\sigma(P)\quad(\sigma\in G_m),\qquad
 B_r\subseteq\Lambda=P B_d,\qquad p\Lambda\subseteq B_r.
\]
For $y\in\mathbb Q_p^d$ the additional integrality assertion is
$Py\in\Lambda\Rightarrow y\in\mathbb Z_p^d$.

If $y$ represents a class of $(V/T)^{G_K}$, then
$M(\sigma)y-y\in\mathbb Z_p^d$ on $G_m\subseteq G_K$.
For $x=pPy$ the differences $U(\sigma)\sigma(x)-x$ therefore
belong to $B_r$. In particular $hx-x\in B_r$ on $H$, so A79L2
writes $p x=z+e$, $z\in K_\infty^r$, $e\in B_r$. Since
$U(g)$ is integral, $(D_g-I)z$ is integral. A79L4 gives
$p^Nz\in B_r$ and hence $p^{N+2}Py\in B_r\subseteq\Lambda$.
The rational integrality assertion gives $p^{N+2}y\in\mathbb Z_p^d$.
For $r=0$ we have $Py=0\in\Lambda$ already and the same conclusion
holds. Enlarging the sum from $J$ to $U_\tau$ now proves
\[
 p^{\,2+\sum_{u\in U_\tau}(m+v_p(u))}(V/T)^{G_K}=0 .
\]
This argument uses a single embedding. A finite union of the sets of
jumps gives the coarser constant in the native theorem: every embedding
jump is a jump of the full filtration, because intersecting two equal
successive filtration steps with the same component preserves equality.
For the interval $[0,1)$ there is only the index zero. The source
adds its nonnegative degree-valuation term and replaces the sum over
nonzero full jumps by its maximum with one. The resulting power is
at least the one just proved.
% END A79L5

% BEGIN A79L6
\paragraph{A79L6. Specializing the integral Wach lattice at the prescribed level.}
We construct the comparison used in A79L5. Take the very $N$ in
$\mathrm{HypW}(T,M_0,h_0)$ and the tied functor and coefficient
inclusion of L78W12--L78W13. In the notation of L78W3 put
$x=[\epsilon_1]-1$ and $H_0=M_0+h_0$. If $j$ is the
given embedding and $A_0=O\otimes_{\mathbb Z_p} A_{\inf}$,
the two height clauses of $N$ give the more precise inclusions
\[
 x^{M_0}j(N)\subseteq A_0\otimes_O T,\qquad
 x^{h_0}(A_0\otimes_O T)\subseteq A_0j(N).
\]
Indeed in the two matrix equations of L78W3 one can use $M_0$
in the first and $h_0$ in the second, since the respective
Frobenius inclusions hold at exactly these powers. Apply L78W2 to
$x^{M_0}Y$ and $x^{h_0}Y^{-1}$ separately. This works for the
basis supplied by $\mathrm{HypW}$, whose image is a basis after
extension to the large Witt ring.

Let $n$ be the least nonnegative integer with $H_0\le p^n(p-1)$,
$j_0=s+n$, and $m=2s+n+1$. Specialize $A_{\inf}$ by
$\theta_n=\theta\circ\varphi^{-n}$, where inverse Frobenius on the
perfect tilt induces the Witt automorphism. This map is equivariant
and sends $[\epsilon_k]$ to $\zeta_{p^{k+n}}$. Extend by the
identity on $O$. The variable $\pi=[\epsilon_s]-1$ becomes
$z=\zeta_{p^{j_0}}-1$, and $x$ becomes
$c=\zeta_{p^{1+n}}-1$. In particular $c\ne0$ and
$c^{H_0}$ divides $p$ in $\mathcal O_C$, since
$v_p(c)=1/(p^n(p-1))$. Define $\Psi(y)$ as the value of
$x^{M_0}j(y)$ in $\mathcal O_C\otimes_{\mathbb Z_p}T$.
The first inclusion makes it integral. The second says that
$c^{H_0}(1\otimes T)$ lies in the
$(\mathcal O_C\otimes_{\mathbb Z_p}O)$-span of $\Psi(N)$.
For $\sigma\in G_m$, equivariance of $j$ and $\theta_n$ gives
$\Psi(\sigma y)=M(\sigma)\sigma(\Psi(y))$ in rational lattice
coordinates: both $x$ and $\sigma x$ specialize to $c$, so the
identities in the big Witt module may be cleared by $x^{M_0}$
on each side and the nonzero scalar $c^{M_0}$ then cancelled.
On $H$ the Wach action on $y$ is the identity.

Choose a $\mathbb Z_p$-basis $o_a$ of $O$ and a series-ring basis
$b_k$ of $N$. Let the columns of $D$ be $\Psi(o_a b_k)$ in the
lattice basis of A79L5. This is square: the basis of $N$ over the
large Witt module has the same cardinality as an $O$-basis of $T$,
and restriction of scalars multiplies both ranks by $[L:\mathbb Q_p]$.
Expanding any series coefficient in the basis $o_a$ after
specialization expresses $\Psi(y)$ as $D$ times an integral
column. The $\mathcal O_C$-span of these columns is already stable
under $O$, since multiplication by $o\in O$ sends $\Psi(y)$ to
$\Psi(oy)$. Thus the previous inclusions yield
\[
          p B_d\subseteq D B_d\subseteq B_d.
\]
In particular $D$ is invertible and $pD^{-1}$ is integral.

The specialized series coefficients have all their $o_a$-coordinates
in $\mathcal O_{K_m}$. For a series in one component, reduce modulo
$\pi^k$ and use the finite expansion at $z$; the remainder tends
to zero since $|z|<1$ and its coordinates before multiplication by
$z^k$ are integral. The scalar-coefficient structure constants in the
$o_a$ basis are in $\mathbb Z_p$. It remains to check the component
idempotents. Over $C$ the algebra $C\otimes_{\mathbb Q_p}L$
splits by the embeddings of $L$. The evaluations of any idempotent
are zero or one, and the embedding matrix of the $o_a$ is invertible
over $\overline{\mathbb Q}_p$ by separability. Its inverse shows that
the coordinates of each specialized idempotent are algebraic.
Equivariance and the fact that the Wach action fixes the series
idempotents show that these coordinates are fixed by $G_m$,
so lie in $K_m$. This proves the coefficient assertion, including
for products of series components.

For $\sigma\in G_m$ express $\sigma(o_a b_k)$ in the Wach basis and
specialize these coefficients. The resulting matrix is necessarily
$U_0(\sigma)=D^{-1}M(\sigma)\sigma(D)$, with entries in
$\mathcal O_{K_m}$. This formula gives continuity, the preceding
field of coefficients gives multiplicativity, and $U_0|_H=I$.
It also shows that scalar multiplication by $a\in O$ in the $D$-basis
has entries in $\mathcal O_{K_m}$, by expanding $a o_b b_k$.

For the smallness at a generator choose $g_0\in G_s$ with
$v_p(\chi(g_0)-1)=s$. Such an automorphism exists because $K_0$
is unramified, the cyclotomic extension is totally ramified,
and $1+p^s$ gives an automorphism of its tower. Put
$g=g_0^{p^{j_0+1}}$. The binomial valuation formula gives
$v_p(\chi(g)-1)=s+j_0+1=m$ and $g\in G_m$.
Set $I=(p,\pi)$ in the series ring and $Y=g_0|_N$.
The congruence clause of $\mathrm{HypW}$ gives $(Y-1)N\subset IN$.
On coefficients $g_0$ fixes constants and the cyclic idempotents,
and sends $\pi$ to a series with constant term zero and linear
coefficient $\chi(g_0)\equiv1\pmod p$, hence fixes $\pi$ modulo $I^2$.
Expanding generators gives $(g_0-1)I^k\subset I^{k+1}$
(in degree zero use the power-series expansion modulo $\pi$).
Semilinearity then gives $(Y-1)(I^kN)\subset I^{k+1}N$.
In $(1+Z)^{p^{j_0+1}}-1$ every coefficient of degree
$0<b<p^{j_0}$ is divisible by $p^2$: use
$b\binom{p^{j_0+1}}{b}=p^{j_0+1}\binom{p^{j_0+1}-1}{b-1}$.
Substituting $Z=Y-1$ proves $(g-1)N\subset(p^2,I^{p^{j_0}})N$.
After specialization all coordinates of this ideal have norm at most
$\max(p^{-2},|z|^{p^{j_0}})<|\delta|$,
since $|p|\le|z|$ and $|z|^{p^{j_0}}=p^{-p/(p-1)}$.
Applied to each basis column this proves $|U_0(g)-I|_{\max}<|\delta|$.
% END A79L6

% BEGIN A79L7
\paragraph{A79L7. The weighted vectors and their embedding component.}
Use the de Rham comparison for the given representation on the tied
Fontaine datum. On the split hub W79H1 gives it from crystallinity.
More generally, in the native $\mathrm{HypW}$ statement,
crystallinity over $K_s$ gives de Rham over $K_s$ by the
crystalline-to-de Rham clause of $\mathrm{fontCore}$ there.
The field $K_s/K$ is finite Galois by goodness and the abelian
cyclotomic construction. Semilinear Galois descent on the space
of periods over $K_s$ preserves its dimension on taking invariants
over $K$, giving the required de Rham comparison over $K$.

Write $t$ for the cyclotomic period, so $\sigma t=\chi(\sigma)t$;
the tied degree-zero residue map $B_{\rm dR}^+\to C$ has kernel
$tB_{\rm dR}^+$. Indeed the tied filtration is the one defined by
the powers of $\xi$, reduction modulo $\xi$ gives Fontaine's
residue map, and $t$ has filtration one with nonzero first symbol.
Thus $t/\xi$ is a unit in the complete local ring and the two
generators give the same filtration. Pick a $K$-basis of the
de Rham period space adapted to its exhaustive separated filtration,
by lifting bases from its finitely many graded pieces. In the rational
lattice coordinates take a basis period $w_k$ of exact filtration $u_k$
and put $e_k=(t^{-u_k}w_k)\bmod t$. These are well-defined
vectors in $C^d$. They satisfy
$\chi(\sigma)^{u_k}M(\sigma)\sigma(e_k)=e_k$ for $\sigma\in G_K$.
Within each weight the $e_k$ are independent over $K$: a relation
would place the corresponding linear combination of the $w_k$ in
the next filtration step. In fact they are independent over $C$.
In a shortest dependence normalize one coefficient to one. Apply
the Galois action and multiply by the character of that selected
vector. Subtracting gives a shorter relation, so each nonzero
coefficient is a scalar of the corresponding character difference.
For a scalar satisfying $\chi(\sigma)^v\sigma(\lambda)=\lambda$,
A79L2 puts it in $\widehat K_\infty$; A79L3 kills its complementary
part. On $K_m$ the generator forces $v=0$ if $\lambda\ne0$,
and invariance under all of $G_K$ then puts $\lambda$ in $K$.
This reduces the proposed dependence to a same-weight relation
over $K$. The de Rham dimension equals $d$, hence these $e_k$
form a basis.

For $\tau:K\to L$, let $E$ project onto the component where
the coefficient action of $\tau(a)$ equals the scalar $a\in K$.
Explicitly apply the coefficient action on the first factor via
$\tau$ and the inclusion into $C$ on the second factor to the
separability idempotent of $K\otimes_{\mathbb Q_p}K$.
This idempotent is one on the identity factor and zero on the others
in the split algebra, so $E$ preserves the de Rham space and its
filtration, commutes with $M(\sigma)$ and is fixed by $G_m$.
The symbol of $Ew_k$ is $Ee_k$; if $u_k$ is not a
$\tau$-jump then the two consecutive filtration steps on that
component coincide, giving $Ee_k=0$. Moreover $D^{-1}ED$ has
entries in $K_m$. This follows by expanding the separability
idempotent: its second-factor scalars lie in $K\subseteq K_m$,
and its first-factor actions are in the $K_m$-span of the integral
coefficient actions computed in A79L6.
% END A79L7

% BEGIN A79L8
\paragraph{A79L8. Integrality after projection and the complete local exponent.}
Put $E'=D^{-1}ED$ over $K_m$. Take $r=\operatorname{rank} E'$.
Choose $r$ rows and $r$ columns whose minor $A'$ has maximal
absolute determinant among all minors of this size. Let $B'$ be
the selected rows, $S$ the inclusion of the selected coordinate
axes, and $R=(A')^{-1}B'$. Cramer's rule gives $R$ integral,
and $S$ is integral. Since $B'S=A'$ and $(E')^2=E'$,
\[
 RS=I,\qquad RE'=R,\qquad E'SR=E'.
\]
The last equality follows by expressing each row of $E'$ in the
chosen row basis and restricting to the selected columns.
The construction includes the zero rank by empty matrices.
Now $E'$ commutes with every $U_0(\sigma)$:
conjugate the commutation with $M(\sigma)$, use the intertwining
in A79L6 and that both $E'$ and $E$ are fixed by $G_m$.
Consequently $R U_0(\sigma) S R=R U_0(\sigma)$, so
$U(\sigma)=R U_0(\sigma)S$ is a representation and
$P=RD^{-1}$ intertwines the two semilinear actions.
All properties of a frame pass to $U$; in particular
$U(g)-I=R(U_0(g)-I)S$. The integral matrix $DS$ is a right
inverse to $P$, and $pP$ is integral. Thus both ball inclusions
of A79L5 hold. Since $P E=P$, the nonzero $P e_k$ have weights
in $U_\tau$ by A79L7. They span $C^r$ by surjectivity; choosing
an independent subset gives the weighted basis of A79L4.

For the last condition on rational coordinates, extend the inverse
of $\tau$ on its image to an embedding $\iota:L\hookrightarrow C$,
so $\iota(\tau(a))=a$. Take an $O$-basis of $T$ and let $Z$ be the
matrix sending rational lattice coordinates to these $L$-coordinates
followed by $\iota$, extended $C$-linearly. All entries of $Z$ are
integral. If $y\in\mathbb Q_p^d$ and $Zy$ is integral, the original
$L$-coordinates are integral since $\iota$ is isometric; hence
$y\in\mathbb Z_p^d$. By the choice of $\iota$, $ZE=Z$.
Also $E=(DE'S)P$, so $Z$ factors through $P$.
If $Py=Pv$ for $v\in B_d$, then $Zy=Zv$ is integral.
This gives exactly the rational test of A79L5 and completes the
construction on the given lattice.

We have therefore proved A79L5's displayed annihilation under
$\mathrm{HypW}$, goodness, the tied period comparisons just specified
and absence of zero for the chosen embedding. On the split hub
W79H1 supplies the de Rham hypothesis. The construction uses neither
a rank-one decomposition of $T$ nor an assumed denominator for
Fontaine's change of basis.
% END A79L8
